\section{Related Work}
\subsection{Zero-Knowledge Proofs for Tree-Based Models}
\label{subsec:related-zk-tree}

Zero-knowledge proofs for tree-based machine learning models have been studied
under several different goals: proving inference, proving accuracy or fairness
properties, and proving training. Early work on private client-side profiling
already considered random forests together with zero-knowledge proofs of correct
model evaluation~\cite{danezis2012private}. Later, Zhang et al.~\cite{zkdt20}
initiated the systematic study of zero-knowledge decision trees. Their zkDT
protocols allow a prover to prove the correctness of a decision-tree prediction
or an accuracy claim without revealing the model. They also extend the approach
to random forests by proving the predictions of individual trees and aggregating
their outputs. This line of work is primarily concerned with inference and
accuracy testing for already trained models, rather than proving that a model is
the output of a prescribed training algorithm.

Several subsequent works improve or extend zero-knowledge proofs for tree
inference. \citet{verifiable-ai22} study verifiable
decentralized AI pipelines and include protocols for verifiable inference of
decision trees and random forests. \citet{cq++24} introduce
lookup argument techniques, including matrix lookup, and apply them to
zero-knowledge decision trees.
They use hyper-box encoding of decision tree,
which, composed with a table-independent lookup argument~\cite{cq22, cq++24},
allows a tree-independent proof of inference for decision tree.

Another related direction is zero-knowledge auditing of fairness for tree-based
models. Confidential-PROFITT~\cite{confidential-profitt22} proposes fair decision
tree learning algorithms together with customized zero-knowledge protocols that
certify fairness while keeping both the model and training data confidential; it
also extends the approach to ensembles of trees. Recent auditing frameworks
further study robustness of fairness certificates and transparent generation of
audit data~\cite{audifair26}. These works prove statistical or fairness properties
of a committed model, or prove that a model was produced by a fairness-aware
training rule. In contrast, our work targets the full training semantics of
gradient-boosted decision trees, including the algebraic constraints for
histogram construction, split selection, leaf-value computation, and boosting
state updates.

The closest works to ours are zero-knowledge proofs of training for forest-based
models. \citet{sparrow24} propose a space-efficient zkSNARK, Sparrow, for
data-parallel circuits and applies it to 
zero-knowledge proofs of random forest training
and prediction.
More
recently, ZKBoost~\cite{zkboost26} presents a zero-knowledge proof of training for
XGBoost. ZKBoost reformulates XGBoost as a deterministic fixed-point algorithm
and gives a generic template instantiated with a general-purpose ZKP
backend. 
Our work is complementary but takes a different approach: rather than
expressing the whole trainer as a general arithmetic circuit, we design
specialized algebraic protocols for the structure of GBDT training. In
particular, we exploit the histogram-based nature of tree construction and batch
the repeated linear, lookup, and aggregation relations that arise across nodes,
features, bins, classes, and boosting rounds.

Overall, prior work on zero-knowledge tree models has mainly focused on
verifiable inference, accuracy/fairness auditing, or general-purpose circuit
encodings of tree training. Our contribution is a specialized proof system for
GBDT training that directly captures the committed training trace and proves
that the committed ensemble is obtained by executing the public
training algorithm on the committed dataset.

\subsection{Batching Lookup Arguments}
As a corollary of \Cref{lem:interleaving},
    we can batch multiple lookup arguments into one,
    even if (1) the contents are of different sizes,
    and (2) the tables are different and potentially of different sizes.
    To be specific,
    given KZG commitments to content vectors $\{\mathbf{f}_i\in \mathbb{F}^{m_i}\}_{i\in [n]}$
    and table vectors $\{\mathbf{t}_i\in \mathbb{F}^{M_i}\}_{i\in [n]}$,
    we want to prove that, \emph{pair-wisely},
    \[
        \mathbf{f}_i \subseteq \mathbf{t}_i\quad \forall i\in [n],
    \]
    where ``$\subseteq$'' denotes that every element in $\mathbf{f}_i$ is also in $\mathbf{t}_i$.

To our best knowledge, this types of batching
    is not studied in the literature.
    Existing works on lookup arguments typically fall into one of the following categories:

\paragraph{Batching many lookups into a common table.}
A common form of batching is to prove that many queried values, or many queried
columns, belong to the same lookup table.
This is the setting handled by standard Plonkish lookup arguments such as
Plookup~\cite{plookup20},
as well as logarithmic-derivative-based lookup arguments such as
LogUp~\cite{logup22}.
These protocols amortize the cost of many membership checks by reducing them to
one multiset relation against a single table.
When multiple tables are supported, the usual approach is to concatenate the
tables and add a table identifier, or tag, so that a lookup into table $i$ is
encoded as a lookup of the tuple $(i,x)$ into one global tagged table.

This technique is effective when all lookups can be represented as membership
in \emph{a single combined table}, but it does not directly batch a family of
\emph{independent} committed relations
\[
    \mathbf{f}_i \subseteq \mathbf{t}_i
    \quad \forall i\in [i],
\]
where both the queried vectors and the corresponding tables are separately
committed and may have unrelated sizes.
We note that committing contents and tables separately is
helpful in our GBDT construction, as it allows us to 
fetch different forest attributes and process them separately.

\paragraph{Large-table lookup arguments.}
Another line of work focuses on reducing the cost of lookup arguments when the
table is large.
Examples include Caulk, Caulk+, flookup, Baloo, cq, and subsequent
improvements~\cite{caulk22,caulk+22,flookup22,baloo25,cq22,locq24}.
The main goal of this line is to make the prover sublinear, table-independent, or
quasi-linear in the number of lookups after preprocessing.
This direction is orthogonal to our batching problem.

\paragraph{Vector, matrix, and segment lookup arguments.}
Recent works also study richer lookup relations, such as vector lookup,
matrix lookup, and segment lookup~\cite{cq++24,sublonk24}.
These arguments can prove that a committed vector, row, subvector, or submatrix
is contained in a larger committed table, often by adding positional labels or
by reducing the relation to a standard lookup.
In our construction, we utilize the matrix lookup of~\cite{cq++24},
which is nevertheless orthogonal to the batching problem.

In summary, prior batching techniques typically batch many queries into one
table, encode multiple tables using tags in one global table, or prove richer
containment relations with respect to one committed table-like object.
By contrast, our interleaving lemma gives a direct way to batch a family of
independent committed lookup relations
$\{(\mathbf{f}_i,\mathbf{t}_i)\}_i$ into 
one single lookup argument,
while allowing both the
query vectors and the tables to have different sizes.
